\documentclass{standalone}
\usepackage{tikz}
\usepackage{graphicx}

\newcommand*\circled[4]{\tikz[baseline=(char.base)]{
    \node[shape=circle, text=#2, fill=#3, draw=#4,
     text opacity=1, fill opacity=1, draw opacity=1,  
     inner sep=2pt, minimum size=0pt] (char) {#1};}}

\makeatletter
\newcommand{\empolygon@impl}[2]{%
 \draw[red, thick, fill=red!20, fill opacity=0] (em#1-1) 
    \foreach \i in {2,...,#2} { -- (em#1-\i) } -- cycle;

 \node (em#1-node) at (em#1-1) [anchor=north west, %xshift=-5pt, yshift=-5pt, 
  draw=yellow, draw opacity=.4, fill=yellow, fill opacity=.1] {\circled{#1}{black}{white}{yellow}};
}

\newcommand{\empolygon}[1]{%
\empolygon@impl{#1}{\@nameuse{em#1-count}}
}
\makeatother


\begin{document}
\begin{tikzpicture}[x=1pt, y=-1pt] 

\node[anchor=south west, inner sep=0, outer sep=0] (image) at (0,0) 
   {\includegraphics[width=%image-width%]{%image-name%}};


\begin{scope}[xscale=1,yscale=1,xshift=0,yshift=%image-height%]

 % pattern is  \coordinate (em#-#) at (x, y);

%edge-map-coordinates%

\end{scope}

\makeatletter
 % pattern is  \@namedef{em#-count}{#}

%edge-map-counts%
\makeatother



 % pattern is  \empolygon{#}
%em-polygons%


\end{tikzpicture}
\end{document}


